% Generated draft: requires mathematical and source review.
\documentclass[11pt]{article}
\usepackage[margin=25mm]{geometry}
\usepackage{fontspec}
\setmainfont{Latin Modern Roman}
\usepackage{amsmath,amssymb,mathtools,unicode-math}
\setmathfont{Latin Modern Math}
\usepackage{enumitem,needspace,xurl,hyperref}
\setlength{\parindent}{0pt}
\setlength{\parskip}{5pt}
\setlength{\emergencystretch}{3em}
\newcommand{\N}{\mathbb{N}}
\newcommand{\Z}{\mathbb{Z}}
\newcommand{\Q}{\mathbb{Q}}
\newcommand{\R}{\mathbb{R}}
\newcommand{\C}{\mathbb{C}}
\newcommand{\F}{\mathbb{F}}
\newcommand{\A}{\mathbb{A}}
\newcommand{\PP}{\mathbb{P}}
\newcommand{\E}{\mathbb{E}}
\newcommand{\eps}{\varepsilon}
\newcommand{\dd}{\,\mathrm{d}}
\begin{document}
\section*{UnsolvedMath ID: 800250}
\textbf{Source problem code:} catalogue entry
\subsection{Definitions and mathematical statement}
For terminology, let \(\Omega\subset\C^n\) be a domain, let \(p\in\partial\Omega\), and suppose a real defining function \(\rho\) of class \(C^m\), with \(m\geq2\), exists near \(p\), so that locally \(\Omega=\{\rho<0\}\) and \(\dd\rho\neq0\) on the boundary. These are explanatory choices; the source does not specify the dimension, boundary regularity, or defining function. The complex tangent space and Levi form are
\[
T_q^{1,0}\partial\Omega=\left\{v\in\C^n:\sum_{j=1}^n\rho_{z_j}(q)v_j=0\right\},\qquad
\mathcal L_\rho(q;v)=\sum_{j,k=1}^n\rho_{z_j\overline z_k}(q)v_j\overline{v_k}.
\]
Here a weakly pseudoconvex point is understood to mean that the Levi form on the complex tangent space is positive semidefinite and degenerate; the source supplies no more precise convention or local geometric hypotheses. Derivatives of order \(r\) of the coefficients of this form require corresponding regularity, for example \(m\geq r+2\), and their interpretation requires choices of coordinates, defining function, and, for tangential components, a tangent frame.

Write \(\Delta=\{\zeta\in\C:|\zeta|<1\}\). One convention for the infinitesimal Kobayashi metric is
\[
k_\Omega(z;v)=\inf\left\{\lambda>0:\text{there exists a holomorphic }f:\Delta\to\Omega,\ f(0)=z,\ \lambda f'(0)=v\right\}.
\]
The source does not specify a metric normalization or nondegeneracy assumptions.

\textbf{Source question.} Can certain derivatives of the Levi form at a weakly pseudoconvex boundary point \(p\) be recovered from the weighted boundary behavior of the Kobayashi metric?

The words ``certain derivatives,'' ``weighted boundary behavior,'' and ``recovered'' are unspecified in the supplied record. In particular, it gives neither the derivative orders and directions, nor the weights, approach regions, asymptotic data, or equivalence under changes of defining function and coordinates. Thus the question is retained as a research target; the supplied evidence does not determine a fully quantified recovery assertion.
\subsection{Short English statement}
The problem asks whether weighted asymptotics of the Kobayashi metric near a weakly pseudoconvex boundary point determine some derivatives of the Levi form there. The derivatives, weights, geometric hypotheses, and meaning of determination remain unspecified, so a precise recovery theorem cannot be formulated from this record alone.
\subsection{Sources}
\begin{itemize}
\item Supplied local UnsolvedMath registry record, permanent ID \(800250\), source code \texttt{catalogue entry}, titled \emph{Several Complex Variables Problem List --- Problem 49}. The registry labels the problem open; this is a supplied status, not an independently verified literature assessment.
\item Buzzard, Fornaess, and Krantz, \emph{The Several Complex Variables Problem List} (1995), Problem 49. Supplied source URL: \url{https://arxiv.org/abs/math/9509202}. The provenance records access on 2026-10-16 and extraction by \texttt{arxiv\_latex}; these are supplied provenance statements, not claims of independent access or verification.
\item The problem text attributes the question to S. Krantz and supplies the contact address \texttt{sk@math.wustl.edu}. The registry lists Buzzard, Fornaess, and Krantz as proposers and 1995 as the proposed year.
\end{itemize}
\end{document}
